\section{Por que é um problema distribuído}

Uma única máquina tem limites de processamento (núcleos de CPU), de memória
e de leitura em disco. Quando o volume de dados cresce, o tempo de
processamento cresce junto e, a partir de certo ponto, os dados nem cabem
na memória de uma máquina só. Distribuir o processamento contorna esses
limites: cada máquina recebe uma parte dos dados, e o trabalho total é
feito em paralelo.

A distribuição, porém, traz desafios que não existem em uma máquina só:
máquinas com desempenho diferente, falhas durante a execução e o custo de
mandar dados pela rede. Esses desafios são o foco dos experimentos
planejados (Seção~\ref{sec:metricas}).
